\documentclass{jsarticle}
\usepackage{amsmath}
\usepackage{amssymb}
\begin{document}
\title{}
\author{}
\date{}
\maketitle

  双子のパラドックスは、一般には、片方が宇宙船に乗り、タイムマシーンで
宇宙の果てまで行き、もう一方が地球に残り、時間が経って戻って来ると、
宇宙船の方の人は、年を取らずに、地球に残った人の方が年を取る。

  私の双子のパラドックスは、片方が、ヒッグス場の密度エネルギーのコントロールで
時の流れの拡散をつくり、宇宙の果てまで、この密度エネルギーのコントロールで、
時間と空間の区別無しに、タイムマシーンで行き、密度エネルギーの拡散で、時間
と空間の物差しがなくなっている状態で地球に戻ると、地球に残った人の時代の前
にも戻れる。この時間の流れの抵抗数では、時が止まれば、空間の物差しがなくなり、
光速や剛体力学とは無縁の時間空間になる。
本当のタイムマシーンになっている。この考えの時間と空間の関係は、
相対性理論を使っている。ヒッグス場の密度エネルギーは抵抗数を表している。
抵抗数が大きいほど、時間の流れの速さは遅くなり、小さいほど、時間の流れは、
自然と流れる。タイムマシーンは、この密度エネルギーの大きさでコントロール
される。年々重力定数は減少しているらしい。ゼータ関数から導かれる
大域的微分方程式が述べる宇宙の終わりには、異次元空間とこの宇宙が
収縮と拡張する空間が一定になり、平坦な宇宙になる。この宇宙は
アインシュタイン博士は、一般相対性理論からの特殊解でもとめていた。
ラムダ項導入は間違いではなかった。
Lisa Randall博士は、knocking on heaven's doorと主題にしていて、AdS5多様体が
平坦な宇宙になったときに、Heavenが５次元多様体と同値になるとも、言っていた。
アカシックレコードのことも、どの時間空間にこの情報空間が接続されているかも、
知っていた。
エーテルは、アインシュタイン博士によって否定されたが、このヒッグス場の
密度エネルギーは、物体の液体とは違う働きをしているのと、実体が宇宙を
覆うブラックホールの密度エネルギーの大きさがその根本にあり、いくら宇宙の
エーテルを探しても、体積と質量とは、密度はちがう。体積が密度になるはずがなく、
質量が密度エネルギーになるはずもなく、体積と質量から生成された、物体の
フェルミオンがこのヒッグス場の密度エネルギーである。目に見えるエーテルを
探しても、意味がない。それ故に、ヒルベルト空間による標数０の体の上の
代数多様体が、大域的微分方程式と同値となる。ゼータ関数の性質の均配が０と
なり、アーベル多様体の収束半径がゼータ関数の閾値となる。この閾値は、原子の
開集合でもあり、量子群は、光量子仮説でもある。このゼータ関数と量子群から
導かれる大域的微分方程式が、ヒッグス場の密度エネルギーである。
特殊相対性理論の慣性質量と重力質量は等価であり、一般相対性理論の
偏曲面による空間の均配が、質量と同じとも言えるのも、空間の性質上、
ヒッグス場の密度エネルギーとは違う。光速度不変の法則は、この密度エネルギー
が一定の場合である。この密度エネルギーは、自然界の慣性力そのままであり、
人為的に今までコントロールされていない。偏曲面による空間の均配の質量と、
異次元空間の均配による質量の斥力が、ヒルベルト空間による標数０の体の上の
多様体を場の理論として、加群としたときの値が、ヒッグス場の密度エネルギー
である。異次元空間がこの宇宙と現存しているようでいないのは、AdS5多様体の
経路積分の式が示している。大域的微分方程式を調べると、積分に初期変数による
多様体を極限値として、その積分を初期関数の変数で代数多様体としている。これは、
慣性質量と重力質量の均配力をも表している。多様体による均配力が大域的微分
方程式でもある。非共変系であるD-braneと慣性力である特殊相対性理論を
一般化したのが、一般共変系であり、この均配力と多様体のフェルミオンを
合わせたのが、ヒッグス場の密度エネルギーである。片方だけでは、密度エネルギー
は求まらない。特殊相対性理論と一般相対性理論と量子力学を合わせた力が、
ヒッグス場の密度エネルギーである。要するに、フェルミオンとボソンを合わせた
超弦理論が、ヒッグス場の密度エネルギーである。これが大域的微分方程式である。

このタイムマシーンの仕組みで、問題になるのが、共時性だ。相対性理論には、
共時性がない。量子力学にはある。タイムマシーンの仕組みで、相対性理論を
使ったと書いた。それは、時間の流れが停止した場合、時間と空間の物差しが
なくなり、空間の移動が、量子テレポーテーションのようになっている。
実際は、意味がちがう。だが、この移動で、宇宙の周りが過去で宇宙の中心
が現在と言えれば、このヒッグス場の密度エネルギーのコントロールで、宇宙船
の行き来が、未来と過去をひっくり返すタイムマシーンの仕組みとなっている。
D-braneの内部でも、相対性理論は成り立ち、D-braneの外側では、assemble-D-brane
として、D-brane同士で相対性理論が成り立つ。生物のDNAには、テロメアが
重力を感知して、生物の寿命をつくり、考古学や物理学でも、地層から年代を
割り出して、相対性原理として、考古学の寿命を形成している。D-braneだけだと、
内部では相対性理論が成り立つが、D-braneの外側では、絶対時空になり、
タイムマシーンの仕組みは使えない。それゆえに、assemble-D-braneだとD-braneの
外側でもタイムマシーンの仕組みは使える。共時性とは、奥が深い考えである。
河合隼雄先生が、ユングとフロイトの本で、この共時性を取り上げていた。
ハイゼンベルクやパウリもこの共時性を研究していた。





\end{document}
